\section{{\v C}ech cohomology as a functor on presheaves}
\label{section-cech-functor}

\noindent
Warning: In this section we work almost exclusively with presheaves and
categories of presheaves and the results are completely wrong in the
setting of sheaves and categories of sheaves!

\medskip\noindent
Let $X$ be a ringed space.
Let $\mathcal{U} : U = \bigcup_{i \in I} U_i$ be an open covering.
Let $\mathcal{F}$ be a presheaf of $\mathcal{O}_X$-modules.
We have the {\v C}ech complex
$\check{\mathcal{C}}^\bullet(\mathcal{U}, \mathcal{F})$
of $\mathcal{F}$ just by thinking of $\mathcal{F}$
as a presheaf of abelian groups. However, each term
$\check{\mathcal{C}}^p(\mathcal{U}, \mathcal{F})$ has a natural
structure of a $\mathcal{O}_X(U)$-module and the differential is given by
$\mathcal{O}_X(U)$-module maps. Moreover, it is clear that the
construction
$$
\mathcal{F} \longmapsto \check{\mathcal{C}}^\bullet(\mathcal{U}, \mathcal{F})
$$
is functorial in $\mathcal{F}$. In fact, it is a functor
\begin{equation}
\label{equation-cech-functor}
\check{\mathcal{C}}^\bullet(\mathcal{U}, -) :
\textit{PMod}(\mathcal{O}_X)
\longrightarrow
\text{Comp}^{+}(\text{Mod}_{\mathcal{O}_X(U)})
\end{equation}
see
Derived Categories, Definition \ref{derived-definition-complexes-notation}
for notation. Recall that the category of bounded below complexes
in an abelian category is an abelian category, see
Homology, Lemma \ref{homology-lemma-cat-cochain-abelian}.

\begin{lemma}
\label{lemma-cech-exact-presheaves}
The functor given by Equation (\ref{equation-cech-functor})
is an exact functor (see Homology, Lemma \ref{homology-lemma-exact-functor}).
\end{lemma}

\begin{proof}
For any open $W \subset U$ the functor
$\mathcal{F} \mapsto \mathcal{F}(W)$ is an additive exact functor
from $\textit{PMod}(\mathcal{O}_X)$ to $\text{Mod}_{\mathcal{O}_X(U)}$.
The terms
$\check{\mathcal{C}}^p(\mathcal{U}, \mathcal{F})$
of the complex are products of these exact functors and hence exact.
Moreover a sequence of complexes is exact if and only if the sequence
of terms in a given degree is exact. Hence the lemma follows.
\end{proof}

\begin{lemma}
\label{lemma-cech-cohomology-delta-functor-presheaves}
Let $X$ be a ringed space.
Let $\mathcal{U} : U = \bigcup_{i \in I} U_i$ be an open covering.
The functors $\mathcal{F} \mapsto \check{H}^n(\mathcal{U}, \mathcal{F})$
form a $\delta$-functor from the abelian category of
presheaves of $\mathcal{O}_X$-modules to the category
of $\mathcal{O}_X(U)$-modules (see
Homology, Definition \ref{homology-definition-cohomological-delta-functor}).
\end{lemma}

\begin{proof}
By
Lemma \ref{lemma-cech-exact-presheaves}
a short exact sequence of presheaves of
$\mathcal{O}_X$-modules
$0 \to \mathcal{F}_1 \to \mathcal{F}_2 \to \mathcal{F}_3 \to 0$
is turned into a short exact sequence of complexes of
$\mathcal{O}_X(U)$-modules. Hence we can use
Homology, Lemma \ref{homology-lemma-long-exact-sequence-cochain}
to get the boundary maps
$\delta_{\mathcal{F}_1 \to \mathcal{F}_2 \to \mathcal{F}_3} :
\check{H}^n(\mathcal{U}, \mathcal{F}_3) \to
\check{H}^{n + 1}(\mathcal{U}, \mathcal{F}_1)$
and a corresponding long exact sequence. We omit the verification
that these maps are compatible with maps between short exact
sequences of presheaves.
\end{proof}


\noindent
In the formulation of the following lemma we use the functor $j_{p!}$ of
extension by $0$ for presheaves of modules
relative to an open immersion $j : U \to X$.
See Sheaves, Section \ref{sheaves-section-open-immersions}. For any open
$W \subset X$ and any presheaf $\mathcal{G}$ of $\mathcal{O}_X|_U$-modules
we have
$$
(j_{p!}\mathcal{G})(W) =
\left\{
\begin{matrix}
\mathcal{G}(W) & \text{if } W \subset U \\
0 & \text{else.}
\end{matrix}
\right.
$$
Moreover, the functor $j_{p!}$ is a left adjoint to the restriction functor
see Sheaves, Lemma \ref{sheaves-lemma-j-shriek-modules}.
In particular we have the following formula
$$
\Hom_{\mathcal{O}_X}(j_{p!}\mathcal{O}_U, \mathcal{F})
=
\Hom_{\mathcal{O}_U}(\mathcal{O}_U, \mathcal{F}|_U)
=
\mathcal{F}(U).
$$
Since the functor $\mathcal{F} \mapsto \mathcal{F}(U)$ is an exact functor
on the category of presheaves we conclude that the presheaf
$j_{p!}\mathcal{O}_U$ is a projective object in the category
$\textit{PMod}(\mathcal{O}_X)$, see
Homology, Lemma \ref{homology-lemma-characterize-projectives}.

\medskip\noindent
Note that if we are given open subsets $U \subset V \subset X$
with associated open immersions $j_U, j_V$, then we have a canonical
map $(j_U)_{p!}\mathcal{O}_U \to (j_V)_{p!}\mathcal{O}_V$. It is the
identity on sections over any open $W \subset U$ and $0$ else.
In terms of the identification
$\Hom_{\mathcal{O}_X}((j_U)_{p!}\mathcal{O}_U, (j_V)_{p!}\mathcal{O}_V) =
(j_V)_{p!}\mathcal{O}_V(U) = \mathcal{O}_V(U)$ it corresponds to
the element $1 \in \mathcal{O}_V(U)$.

\begin{lemma}
\label{lemma-cech-map-into}
Let $X$ be a ringed space.
Let $\mathcal{U} : U = \bigcup_{i \in I} U_i$ be an open covering.
Denote $j_{i_0\ldots i_p} : U_{i_0 \ldots i_p} \to X$ the open immersion.
Consider the chain complex $K(\mathcal{U})_\bullet$
of presheaves of $\mathcal{O}_X$-modules
$$
\ldots
\to
\bigoplus_{i_0i_1i_2} (j_{i_0i_1i_2})_{p!}\mathcal{O}_{U_{i_0i_1i_2}}
\to
\bigoplus_{i_0i_1} (j_{i_0i_1})_{p!}\mathcal{O}_{U_{i_0i_1}}
\to
\bigoplus_{i_0} (j_{i_0})_{p!}\mathcal{O}_{U_{i_0}}
\to 0 \to \ldots
$$
where the last nonzero term is placed in degree $0$
and where the map
$$
(j_{i_0\ldots i_{p + 1}})_{p!}\mathcal{O}_{U_{i_0\ldots i_{p + 1}}}
\longrightarrow
(j_{i_0\ldots \hat i_j \ldots i_{p + 1}})_{p!}
\mathcal{O}_{U_{i_0\ldots \hat i_j \ldots i_{p + 1}}}
$$
is given by $(-1)^j$ times the canonical map.
Then there is an isomorphism
$$
\Hom_{\mathcal{O}_X}(K(\mathcal{U})_\bullet, \mathcal{F})
=
\check{\mathcal{C}}^\bullet(\mathcal{U}, \mathcal{F})
$$
functorial in $\mathcal{F} \in \Ob(\textit{PMod}(\mathcal{O}_X))$.
\end{lemma}

\begin{proof}
We saw in the discussion just above the lemma that
$$
\Hom_{\mathcal{O}_X}(
(j_{i_0\ldots i_p})_{p!}\mathcal{O}_{U_{i_0\ldots i_p}},
\mathcal{F})
=
\mathcal{F}(U_{i_0\ldots i_p}).
$$
Hence we see that it is indeed the case that the direct sum
$$
\bigoplus\nolimits_{i_0 \ldots i_p}
(j_{i_0 \ldots i_p})_{p!}\mathcal{O}_{U_{i_0 \ldots i_p}}
$$
represents the functor
$$
\mathcal{F}
\longmapsto
\prod\nolimits_{i_0\ldots i_p} \mathcal{F}(U_{i_0\ldots i_p}).
$$
Hence by Categories, Yoneda Lemma \ref{categories-lemma-yoneda}
we see that there is a complex $K(\mathcal{U})_\bullet$ with terms
as given. It is a simple matter to see that the maps are as given
in the lemma.
\end{proof}

\begin{lemma}
\label{lemma-homology-complex}
Let $X$ be a ringed space.
Let $\mathcal{U} : U = \bigcup_{i \in I} U_i$ be an open covering.
Let $\mathcal{O}_\mathcal{U} \subset \mathcal{O}_X$
be the image presheaf of the map
$\bigoplus j_{p!}\mathcal{O}_{U_i} \to \mathcal{O}_X$.
The chain complex $K(\mathcal{U})_\bullet$ of presheaves
of Lemma \ref{lemma-cech-map-into} above has homology presheaves
$$
H_i(K(\mathcal{U})_\bullet) =
\left\{
\begin{matrix}
0 & \text{if} & i \not = 0 \\
\mathcal{O}_\mathcal{U} & \text{if} & i = 0
\end{matrix}
\right.
$$
\end{lemma}

\begin{proof}
Consider the extended complex $K^{ext}_\bullet$ one gets by putting
$\mathcal{O}_\mathcal{U}$ in degree $-1$ with the obvious map
$K(\mathcal{U})_0 =
\bigoplus_{i_0} (j_{i_0})_{p!}\mathcal{O}_{U_{i_0}} \to
\mathcal{O}_\mathcal{U}$.
It suffices to show that taking sections of this extended complex over
any open $W \subset X$ leads to an acyclic complex.
In fact, we claim that for every $W \subset X$ the complex
$K^{ext}_\bullet(W)$ is homotopy equivalent to the zero complex.
Write $I = I_1 \amalg I_2$ where $W \subset U_i$ if and only
if $i \in I_1$.

\medskip\noindent
If $I_1 = \emptyset$, then the complex $K^{ext}_\bullet(W) = 0$ so there is
nothing to prove.

\medskip\noindent
If $I_1 \not = \emptyset$, then
$\mathcal{O}_\mathcal{U}(W) = \mathcal{O}_X(W)$
and
$$
K^{ext}_p(W) =
\bigoplus\nolimits_{i_0 \ldots i_p \in I_1} \mathcal{O}_X(W).
$$
This is true because of the simple description of the presheaves
$(j_{i_0 \ldots i_p})_{p!}\mathcal{O}_{U_{i_0 \ldots i_p}}$.
Moreover, the differential of the complex $K^{ext}_\bullet(W)$
is given by
$$
d(s)_{i_0 \ldots i_p} =
\sum\nolimits_{j = 0, \ldots, p + 1} \sum\nolimits_{i \in I_1}
(-1)^j s_{i_0 \ldots i_{j - 1} i i_j \ldots i_p}.
$$
The sum is finite as the element $s$ has finite support.
Fix an element $i_{\text{fix}} \in I_1$. Define a map
$$
h : K^{ext}_p(W) \longrightarrow K^{ext}_{p + 1}(W)
$$
by the rule
$$
h(s)_{i_0 \ldots i_{p + 1}} =
\left\{
\begin{matrix}
0 & \text{if} & i_0 \not = i_{\text{fix}} \\
s_{i_1 \ldots i_{p + 1}} & \text{if} & i_0 = i_{\text{fix}}
\end{matrix}
\right.
$$
We will use the shorthand
$h(s)_{i_0 \ldots i_{p + 1}} = (i_0 = i_{\text{fix}}) s_{i_1 \ldots i_p}$
for this. Then we compute
\begin{eqnarray*}
& & (dh + hd)(s)_{i_0 \ldots i_p} \\
& = &
\sum_j \sum_{i \in I_1} (-1)^j h(s)_{i_0 \ldots i_{j - 1} i i_j \ldots i_p}
+
(i = i_0) d(s)_{i_1 \ldots i_p} \\
& = &
s_{i_0 \ldots i_p} +
\sum_{j \geq 1}\sum_{i \in I_1}
(-1)^j (i_0 = i_{\text{fix}}) s_{i_1 \ldots i_{j - 1} i i_j \ldots i_p}
+
(i_0 = i_{\text{fix}}) d(s)_{i_1 \ldots i_p}
\end{eqnarray*}
which is equal to $s_{i_0 \ldots i_p}$ as desired.
\end{proof}

\begin{lemma}
\label{lemma-cech-cohomology-derived-presheaves}
Let $X$ be a ringed space.
Let $\mathcal{U} : U = \bigcup_{i \in I} U_i$
be an open covering of $U \subset X$.
The {\v C}ech cohomology functors $\check{H}^p(\mathcal{U}, -)$
are canonically isomorphic as a $\delta$-functor to
the right derived functors of the functor
$$
\check{H}^0(\mathcal{U}, -) :
\textit{PMod}(\mathcal{O}_X)
\longrightarrow
\text{Mod}_{\mathcal{O}_X(U)}.
$$
Moreover, there is a functorial quasi-isomorphism
$$
\check{\mathcal{C}}^\bullet(\mathcal{U}, \mathcal{F})
\longrightarrow
R\check{H}^0(\mathcal{U}, \mathcal{F})
$$
where the right hand side indicates the right derived functor
$$
R\check{H}^0(\mathcal{U}, -) :
D^{+}(\textit{PMod}(\mathcal{O}_X))
\longrightarrow
D^{+}(\mathcal{O}_X(U))
$$
of the left exact functor $\check{H}^0(\mathcal{U}, -)$.
\end{lemma}

\begin{proof}
Note that the category of presheaves of $\mathcal{O}_X$-modules
has enough injectives, see
Injectives, Proposition \ref{injectives-proposition-presheaves-modules}.
Note that $\check{H}^0(\mathcal{U}, -)$ is a left exact functor
from the category of presheaves of $\mathcal{O}_X$-modules
to the category of $\mathcal{O}_X(U)$-modules.
Hence the derived functor and the right derived functor exist, see
Derived Categories, Section \ref{derived-section-right-derived-functor}.

\medskip\noindent
Let $\mathcal{I}$ be an injective presheaf of $\mathcal{O}_X$-modules.
In this case the functor $\Hom_{\mathcal{O}_X}(-, \mathcal{I})$
is exact on $\textit{PMod}(\mathcal{O}_X)$. By
Lemma \ref{lemma-cech-map-into} we have
$$
\Hom_{\mathcal{O}_X}(K(\mathcal{U})_\bullet, \mathcal{I})
=
\check{\mathcal{C}}^\bullet(\mathcal{U}, \mathcal{I}).
$$
By Lemma \ref{lemma-homology-complex} we have that $K(\mathcal{U})_\bullet$ is
quasi-isomorphic to $\mathcal{O}_\mathcal{U}[0]$. Hence by
the exactness of Hom into $\mathcal{I}$ mentioned above we see
that $\check{H}^i(\mathcal{U}, \mathcal{I}) = 0$ for all
$i > 0$. Thus the $\delta$-functor $(\check{H}^n, \delta)$
(see Lemma \ref{lemma-cech-cohomology-delta-functor-presheaves})
satisfies the assumptions of
Homology, Lemma \ref{homology-lemma-efface-implies-universal},
and hence is a universal $\delta$-functor.

\medskip\noindent
By
Derived Categories, Lemma \ref{derived-lemma-higher-derived-functors}
also the sequence $R^i\check{H}^0(\mathcal{U}, -)$
forms a universal $\delta$-functor. By the uniqueness of universal
$\delta$-functors, see
Homology, Lemma \ref{homology-lemma-uniqueness-universal-delta-functor}
we conclude that
$R^i\check{H}^0(\mathcal{U}, -) = \check{H}^i(\mathcal{U}, -)$.
This is enough for most applications
and the reader is suggested to skip the rest of the proof.

\medskip\noindent
Let $\mathcal{F}$ be any presheaf of $\mathcal{O}_X$-modules.
Choose an injective resolution $\mathcal{F} \to \mathcal{I}^\bullet$
in the category $\textit{PMod}(\mathcal{O}_X)$.
Consider the double complex
$\check{\mathcal{C}}^\bullet(\mathcal{U}, \mathcal{I}^\bullet)$ with terms
$\check{\mathcal{C}}^p(\mathcal{U}, \mathcal{I}^q)$.
Consider the associated total complex
$\text{Tot}(\check{\mathcal{C}}^\bullet(\mathcal{U}, \mathcal{I}^\bullet))$,
see Homology, Definition \ref{homology-definition-associated-simple-complex}.
There is a map of complexes
$$
\check{\mathcal{C}}^\bullet(\mathcal{U}, \mathcal{F})
\longrightarrow
\text{Tot}(\check{\mathcal{C}}^\bullet(\mathcal{U}, \mathcal{I}^\bullet))
$$
coming from the maps
$\check{\mathcal{C}}^p(\mathcal{U}, \mathcal{F})
\to \check{\mathcal{C}}^p(\mathcal{U}, \mathcal{I}^0)$
and there is a map of complexes
$$
\check{H}^0(\mathcal{U}, \mathcal{I}^\bullet)
\longrightarrow
\text{Tot}(\check{\mathcal{C}}^\bullet(\mathcal{U}, \mathcal{I}^\bullet))
$$
coming from the maps
$\check{H}^0(\mathcal{U}, \mathcal{I}^q) \to
\check{\mathcal{C}}^0(\mathcal{U}, \mathcal{I}^q)$.
Both of these maps are quasi-isomorphisms by an application of
Homology, Lemma \ref{homology-lemma-double-complex-gives-resolution}.
Namely, the columns of the double complex are exact in positive degrees
because the {\v C}ech complex as a functor is exact
(Lemma \ref{lemma-cech-exact-presheaves})
and the rows of the double complex are exact in positive degrees
since as we just saw the higher {\v C}ech cohomology groups of the injective
presheaves $\mathcal{I}^q$ are zero.
Since quasi-isomorphisms become invertible
in $D^{+}(\mathcal{O}_X(U))$ this gives the last displayed morphism
of the lemma. We omit the verification that this morphism is
functorial.
\end{proof}





